% !TeX root = ../../Book5.tex
% 纲要标题：### 12.2 Lie 定理
% 纲要位置：工作记录/计划纲要.md，第 4459 行。

\section{Lie 定理}
\label{b5:sec:1202}

% 纲要范围：
%
% * 可解复 Lie 代数的共同特征向量
% * 表示的同时上三角化
% * 特征零与代数闭假设的作用
% * 用仿射自然模落实余维一理想的归纳，展开循环空间取迹，并证明正特征 Heisenberg 表示不可约
%

严格上三角矩阵的对角线全为零，而一般上三角矩阵允许非零特征值。相应地，可解 Lie 代数的表示通常没有公共零向量，却在复数域上总有公共特征向量。证明需要先说明一个理想的特征空间何时被整个 Lie 代数保持。

\subsection{可解复 Lie 代数的共同特征向量}
\label{b5:subsec:1202:01}

\begin{lemma}[理想的特征空间]\label{b5:lem:ideal-weight-stable}
	设 \(\mathfrak g\) 是有限维复 Lie 代数，\(V\) 为有限维复向量空间，\(\rho:\mathfrak g\rightarrow\mathfrak{gl}(V)\) 是表示，\(\mathfrak a\subseteq\mathfrak g\) 是理想，\(\lambda:\mathfrak a\rightarrow\mathbb C\) 是线性泛函。若
	\[
		V_\lambda=\bigl\{\,v\in V\mid \rho(a)v=\lambda(a)v \;\text{对所有}\;a\in\mathfrak a\,\bigr\}\ne0,
	\]
	则 \(\lambda\bigl([\mathfrak g,\mathfrak a]\bigr)=0\)，且 \(V_\lambda\) 是 \(\mathfrak g\)-子模。
\end{lemma}
\begin{proof}
	省略作用符号 \(\rho\)。固定 \(x\in\mathfrak g\) 与非零 \(v\in V_\lambda\)。 对 \(a\in\mathfrak a\)，先把前两层的作用写出：
	\[
		\begin{aligned}
			av&=\lambda(a)v,\\
			a(xv)&=x(av)+[a,x]v =\lambda(a)xv+\lambda([a,x])v.
		\end{aligned}
	\]
	若 \(v,xv\) 线性无关，\(a\) 在它们张成的空间上具有矩阵
	\[
		\begin{pmatrix}
			\lambda(a)&\lambda([a,x])\\
			0&\lambda(a)
		\end{pmatrix}.
	\]
	但这个空间未必被 \(x\) 保持，因而还不能在其中把 \([a,x]\) 当作两个限制算子的交换子取迹。若 \(v,xv\) 线性相关，则 \(\mathbb Cv\) 已被 \(x\) 保持，所需循环空间只有一维。统一处理这两种情况，令 \(m\ge1\) 是满足
	\[
		x^mv\in\operatorname{span}_{\mathbb C}\{v,xv,\ldots,x^{m-1}v\}
	\]
	的最小整数；有限维性保证这样的 \(m\) 存在。于是 \(v,xv,\ldots,x^{m-1}v\) 线性无关，构成循环空间
	\[
		W=\operatorname{span}_{\mathbb C}\{v,xv,\ldots,x^{m-1}v\}
	\]
	的一组基，并且 \(xW\subseteq W\)。对 \(j\ge1\) 归纳，由
	\[
		a x^jv=x a x^{j-1}v+[a,x]x^{j-1}v,\qquad[a,x]\in\mathfrak a,
	\]
	得到
	\[
		a x^jv-\lambda(a)x^jv \in\operatorname{span}_{\mathbb C}\{v,\ldots,x^{j-1}v\}.
	\]
	归纳时第二项对 \([a,x]\in\mathfrak a\) 使用同一归纳断言；第一项中较低次幂经 \(x\) 后至多升到 \(x^{j-1}v\)。因此 \(\mathfrak a\) 保持 \(W\)，每个 \(a\in\mathfrak a\) 在这组基下都上三角，对角元全为 \(\lambda(a)\)，所以 \(\operatorname{tr}_W a=m\lambda(a)\)。这也适用于 \([a,x]\in\mathfrak a\)。此时 \(a,x\) 都保持 \(W\)，记 \(A=a|_W\)、\(X=x|_W\)，就有
	\[
		m\lambda([a,x])=\operatorname{tr}_W[a,x] =\operatorname{tr}(AX-XA)=0.
	\]
	复数域上 \(m\ne0\)，故 \(\lambda\bigl([a,x]\bigr)=0\)。最后，对任意 \(w\in V_\lambda\) 有
	\[
		a(xw)=x(aw)+[a,x]w=\lambda(a)xw,
	\]
	所以 \(xw\in V_\lambda\)。
\end{proof}

\begin{theorem}[Lie 定理的特征向量形式]\label{b5:thm:lie-common-eigenvector}
	设 \(\mathfrak g\) 是有限维可解复 Lie 代数，\(V\ne0\) 是有限维复表示。则 \(V\) 含有一维 \(\mathfrak g\)-子模，即存在 \(v\ne0\) 和线性泛函 \(\lambda:\mathfrak g\rightarrow\mathbb C\)，使 \(xv=\lambda(x)v\) 对所有 \(x\in\mathfrak g\) 成立。
\end{theorem}
\begin{proof}
	对 \(\dim\mathfrak g\) 归纳。零代数时任取非零向量。非零可解代数满足 \([\mathfrak g,\mathfrak g]\ne\mathfrak g\)，否则导出列不能终止。取交换商 \(\mathfrak g/[\mathfrak g,\mathfrak g]\) 中一个余维一子空间，其逆像 \(\mathfrak a\) 是余维一理想。它仍可解，故归纳假设给出非零的 \(\mathfrak a\)-公共特征空间 \(V_\mu\)。由\cref{b5:lem:ideal-weight-stable}，这个空间被 \(\mathfrak g\) 保持。

	取 \(x\in\mathfrak g\setminus\mathfrak a\)，则有向量空间直和分解 \(\mathfrak g=\mathfrak a\oplus\mathbb Cx\)。由于复数域代数闭，\(x|_{V_\mu}\) 有特征向量 \(v\ne0\)，记 \(xv=cv\)。定义线性泛函
	\[
		\lambda:\mathfrak g\longrightarrow\mathbb C,\qquad \lambda(a+tx)=\mu(a)+tc \quad(a\in\mathfrak a,\ t\in\mathbb C).
	\]
	由 \(v\in V_\mu\) 可得 \((a+tx)v=(\mu(a)+tc)v=\lambda(a+tx)v\)，故 \(v\) 是 \(\mathfrak g\) 的公共特征向量。
\end{proof}

在\secref{b5:sec:1103}的仿射自然模中，上述归纳可以直接算出。取复 Lie 代数
\[
	\mathfrak g=\mathbb Ch\oplus\mathbb Ce,\qquad[h,e]=e,
\]
在 \(V=\mathbb C^2\) 上令 \(\rho(h)=E_{11}\)、\(\rho(e)=E_{12}\)，标准基记为 \(e_1,e_2\)。因为 \([\mathfrak g,\mathfrak g]=\mathbb Ce\)，可以选余维一理想 \(\mathfrak a=\mathbb Ce\)。它的公共特征空间为
\[
	V_\mu=\ker\rho(e)=\mathbb Ce_1,\qquad\mu(e)=0.
\]
剩下的一个算子是 \(\rho(h)\)；它保持这条直线，且 \(\rho(h)e_1=e_1\)。因此 \(e_1\) 是整个 \(\mathfrak g\) 的公共特征向量，其特征泛函满足 \(\lambda(h)=1\)、\(\lambda(e)=0\)。这里先处理理想，再在其特征空间上处理最后一个算子，正是证明中归纳的两步。相反，\(\mathbb Ch\) 虽也是余维一子代数，却不是理想；它的零特征空间 \(\mathbb Ce_2\) 不被 \(\rho(e)\) 保持，因为 \(\rho(e)e_2=e_1\)。理想条件保证的正是继续这一步所需的不变性。

\subsection{表示的同时上三角化}
\label{b5:subsec:1202:02}

\term[Lie-dingli]{Lie 定理}[Lie's theorem]通常表述为下面的形式。

\begin{theorem}[Lie 三角化定理]\label{b5:thm:lie-triangularization}
	有限维可解复 Lie 代数 \(\mathfrak g\) 的任意有限维复表示 \(\rho:\mathfrak g\rightarrow\mathfrak{gl}(V)\) 都可同时上三角化。等价地，\(V\) 有一条每个相邻商都为一维的 \(\mathfrak g\)-子模旗标。
\end{theorem}
\begin{proof}
	\(V=0\) 时结论直接成立。否则，\cref{b5:thm:lie-common-eigenvector}给出一维子模 \(V_1\)。在 \(V/V_1\) 上再用同一定理，并提升所得旗标。对 \(\dim V\) 归纳即得完全旗标。选取与之适应的基后，全部作用矩阵均上三角。
\end{proof}

\begin{proposition}\label{b5:prop:solvable-derived-nilpotent}
	设 \(\mathfrak g\) 是有限维可解复 Lie 代数，\(V\) 为有限维复向量空间，\(\rho:\mathfrak g\to\mathfrak{gl}(V)\) 是表示。则 \(\rho\bigl([\mathfrak g,\mathfrak g]\bigr)\) 中的每个算子幂零，而且
	\[
		\operatorname{tr}\bigl(\rho(x)\rho(y)\bigr)=0 \quad\bigl(x\in\mathfrak g,\ y\in[\mathfrak g,\mathfrak g]\bigr).
	\]
	此外，\([\mathfrak g,\mathfrak g]\) 本身是幂零 Lie 代数，\(\mathfrak g\) 的每个非零不可约有限维复表示都是一维的。
\end{proposition}
\begin{proof}
	在\cref{b5:thm:lie-triangularization}给出的基下，上三角矩阵的交换子对角元全为零，所以 \(\rho(y)\) 严格上三角，\(\rho(x)\rho(y)\) 也严格上三角。前两个结论随之成立。取 \(\rho=\operatorname{ad}_{\mathfrak g}\)，再限制到不变子空间 \([\mathfrak g,\mathfrak g]\)，由\cref{b5:thm:engel}得到导出代数幂零。最后，不可约表示中的一维不变子空间只能是整个空间。
\end{proof}

这条命题不能把原代数的可解性加强为幂零性。仍取刚才的仿射代数 \(\mathfrak g=\mathbb Ch\oplus\mathbb Ce\)、\([h,e]=e\)，导出代数 \(\mathbb Ce\) 内部的括号全为零，因而它幂零；但原代数在这个理想上的作用满足
\[
	[\mathbb Ce,\mathbb Ce]=0,\qquad [\mathfrak g,\mathbb Ce]=\mathbb Ce.
\]
在下降中心列中，后一个等式会反复出现，因而 \(\mathfrak g\) 不幂零。导出代数自身的内部括号终止，与原代数反复作用于导出代数，是两个不同的条件。

\ProseParagraphBreak

三角化不等于对角化。令一维交换 Lie 代数的生成元作用为非零 Jordan 幂零块，就得到不可分解而非不可约的表示；它的完全旗标存在，但通常没有不变直和补。

\subsection{特征零与代数闭假设的作用}
\label{b5:subsec:1202:03}

\begin{example}\label{b5:exm:lie-real-failure}
	实一维交换 Lie 代数 \(\mathbb Rx\) 在 \(\mathbb R^2\) 上作用为
	\[
		\rho(x)=\begin{pmatrix}0&-1\\1&0\end{pmatrix}.
	\]
	特征多项式 \(T^2+1\) 在 \(\mathbb R\) 中没有根，因此不存在实不变直线。复化后则分成特征值 \(i\)、\(-i\) 对应的两条直线。
\end{example}

\begin{example}\label{b5:exm:lie-positive-characteristic}
	设 \(k\) 是特征为素数 \(p\) 的域，\(V=k[t]/(t^p)\)。乘以 \(t\) 的算子记为 \(T\)，形式求导记为 \(D\)；对任意 \(f\in k[t]\)，有 \(D(t^pf)=t^pD(f)\)，所以求导在商空间上良定义，且
	\[
		[D,T]=I,\qquad[D,I]=[T,I]=0.
	\]
	三个算子线性无关：若 \(aD+bT+cI=0\)，先作用于 \(1\) 得 \(b=c=0\)，再作用于 \(t\) 得 \(a=0\)。故 \(L=\operatorname{span}_k\{D,T,I\}\) 是三维 Lie 代数，其导出代数为 \(kI\)，且 \([L,kI]=0\)，所以它是二步幂零的。然而 \(D,T\) 不可能有公共特征向量：若 \(v\ne0\)、\(Dv=av\)、\(Tv=bv\)，则 \([D,T]v=(ab-ba)v=0\)，与 \([D,T]v=v\ne0\) 矛盾。 在 \(p=2\) 时，基 \(1,t\) 下的矩阵为
	\[
		\begin{aligned}
			D&=\begin{pmatrix}0&1\\0&0\end{pmatrix},\qquad
			T=\begin{pmatrix}0&0\\1&0\end{pmatrix},
			\\
			DT-TD&=\begin{pmatrix}1&0\\0&-1\end{pmatrix}=I.
		\end{aligned}
	\]
	虽然 \(D^2=T^2=0\)，交换子却是单位算子。此时 \(\operatorname{tr}I=2=0\) 于 \(k\)，所以 \([D,T]=I\) 与交换子的迹为零相容。在任意特征零的非零有限维空间中，同一关系则会迫使 \(0=\operatorname{tr}I=\dim V\)，不可能成立。
\end{example}

\begin{proposition}[正特征 Heisenberg 表示的不可约性]\label{b5:prop:positive-heisenberg-irreducible}
	设 \(k\) 是特征为素数 \(p\) 的域，\(\mathfrak h=ku\oplus kv\oplus kz\) 是由
	\[
		[u,v]=z,\qquad[u,z]=[v,z]=0
	\]
	确定的三维 Heisenberg Lie 代数。令 \(V=k[t]/(t^p)\)，并令 \(D\)、\(T\)、\(I\) 分别为形式求导、乘以 \(t\) 与恒等算子。则
	\[
		\rho(u)=D,\qquad\rho(v)=T,\qquad\rho(z)=I
	\]
	定义了 \(\mathfrak h\) 的一个 \(p\) 维不可约表示。
\end{proposition}
\begin{proof}
	\cref{b5:exm:lie-positive-characteristic}的计算说明 \(D,T\) 在商空间上良定义，并满足表示所需的三个括号关系。设 \(0\ne U\subseteq V\) 是子模。取非零元素的唯一代表
	\[
		f=a_0+a_1t+\cdots+a_rt^r,\qquad 0\le r<p,\quad a_r\ne0.
	\]
	由于 \(U\) 被 \(D\) 保持，\(D^rf=r!a_r\in U\)。这里 \(r<p\)，所以 \(r!\ne0\) 于 \(k\)，从而 \(1\in U\)。再利用 \(T\) 的稳定性，得到
	\[
		T^j1=t^j\in U\qquad(0\le j<p).
	\]
	这些向量构成 \(V\) 的一组基，故 \(U=V\)，证明了不可约性。
\end{proof}

因此，正特征下的幂零 Lie 代数确实可以有维数大于一的不可约表示。这里 \(I\) 不是幂零算子，所以不满足 Engel 定理算子形式的假设；幂零的是抽象 Lie 代数的括号结构。

\ProseParagraphBreak

另一方面，\cref{b5:lem:ideal-weight-stable}中消去 \(m\) 的步骤要求它于基域中非零。在本例中，取理想 \(\mathfrak a=kD\oplus kI\) 与向量 \(1\in V\)，其特征泛函满足 \(\lambda(D)=0\)、\(\lambda(I)=1\)。以 \(x=T\) 生成的循环空间就是整个 \(V\)，维数 \(m=p\)，因而取迹只给出
\[
	m\lambda([D,T])=p\lambda(I)=p\cdot1=0.
\]
它无法迫使 \(\lambda(I)\) 为零；实际 \(T1=t\) 已不在原特征空间中，因为 \(Dt=1\ne0\)。正特征在这里使理想特征空间的不变性失效。
